% !PRINTED 15
\weberpage{15}

%% [running head, not body text -- commented out 2026-08-07] Page furniture (running
%% head, folio, signature) is excluded from this edition by standing policy.
%% \gesperrt{\S.\,5.} \hfill \gesperrt{Mehrstufiger Isomorphismus.} \hfill 15

%<<CONT
bezeichnen sie nach dem Vorgange von Hölder\footnote{Mathematische Annalen, Bd.\,34. Das Zeichen erinnert an einen Quotienten, mit dem ja die complementäre Gruppe eine gewisse Analogie hat.} mit $P/N$. Der Grad der complementären Gruppe ist gleich dem Index der Gruppe $N$.

Wir erwähnen am Schluss dieser Betrachtungen noch einen Satz über Normaltheiler.

\noindent 5.\ \gesperrt{Ist $N$ ein Theiler von $P$, $N'$ ein Theiler von $N$ und zugleich Normaltheiler von $P$, so ist auch $N'$ Normaltheiler von $N$.}

Dies ist selbstverständlich, weil für jedes Element $b$ von $P$ $b^{-1}N'b = N'$ ist, und weil jedes Element von $N$ auch ein Element von $P$ ist.

\gesperrt{Man darf aber nicht umgekehrt schliessen, dass, wenn $N$ ein Normaltheiler von $P$, $N'$ ein Normaltheiler von $N$ ist, auch $N'$ ein Normaltheiler von $P$ sein müsse.}

\begin{center}
\S.\,5.

\gesperrt{Mehrstufiger Isomorphismus.}
\end{center}

Ist wie oben $N$ ein Normaltheiler von $P$ vom Grade $\nu$ und Index $\mu$, also $n = \mu \nu$, so ist die zu $P$ complementäre Gruppe $P/N$ vom Grade $\mu$. Wir wollen diese oder eine damit isomorphe Gruppe mit $Q$ bezeichnen und ihre Elemente, die den Nebengruppen $N, Nb_1, Nb_2 \dots$ entsprechen, mit $A, B, C \dots$ Jedem dieser Elemente entsprechen $\nu$ Elemente der Gruppe $P$, etwa

\noindent dem Elemente $A$ die Elemente $a, a_1 \dots a_{\nu-1}$

\noindent dem Elemente $B$ die Elemente $b, b_1 \dots b_{\nu-1}$

\noindent $. \; . \; . \; . \; . \; . \; . \; . \; . \; . \; . \; . \; . \; . \; . \; . \; . \; . \; . \; . \; .$

Dies Entsprechen ist dann derart, dass jedes zusammengesetzte Element $ab$ dem zusammengesetzten Elemente $AB$ entspricht. Denn die Elemente $A$ und $B$ entsprechen den Nebengruppen $Na$ und $Nb$, und es ist, weil $N$ ein Normaltheiler ist,
\[
Na \, Nb = Nab.
\]

Es gilt also hier bei der Zusammensetzung der $A, B \dots$ einerseits und der $a, b \dots$ andererseits dasselbe Gesetz, wie bei
%CONT>>
